\chapter{Reflection}
The goal of fairness-aware AI is to reduce unfair outcomes caused by automated decision-making. This thesis investigates how fairness constraints can be applied more effectively during model training, with the aim of improving the development of AI systems that satisfy predefined fairness criteria. However, achieving fairness according to a mathematical definition does not necessarily mean that an AI system is fair in a broader societal sense.

This reflection considers the relationship between technical fairness and societal fairness. It discusses the contribution and limitations of fairness constraints, the question of who should determine what constitutes fairness, and whether making an AI model fair can address the underlying causes of unfair outcomes. The purpose is not to question the value of fairness-aware AI, but to consider how such techniques can contribute to fairer systems while recognising that AI remains a tool within a broader social context.

\section{Contributions}
The research in this thesis contributes to a better understanding of how fairness constraints behave during model training. Rather than only evaluating the fairness of the final model, the experiments investigate how the timing of a constraint
%, the definition of the fairness violation,%
and the combination of multiple fairness constraints affect the training process and the resulting models. This provides additional insight into the practical use of fairness-aware training methods. 

A better understanding of these training behaviours can contribute to the development of more effective fairness interventions. For example, the results suggest that fairness constraints do not necessarily need to be applied throughout the entire training process, which can reduce the amount of required computations, or allow more compute to be spent on another part, during the development of an AI model. Similarly, understanding which fairness definitions can be combined and how their strengths interact can help practitioners select more appropriate configurations for a given application.

The societal value of these contributions lies primarily in providing better tools in fair AI. However, these improvements rely on the proper usage of said tools, they only address the fairness properties that are explicitly measured and optimised. They therefore improve the ability to implement a chosen notion of fairness, rather than determining whether that notion is appropriate in a particular societal context.

\section{Who determines fairness?}
An important question in the development of fair machine learning systems is who decides what fairness means. Fairness is not a purely technical property that can be determined independently of the context in which an AI system is used. Instead, developers must choose a fairness definition and a corresponding fairness metric, and this choice can have significant consequences for the people affected by the system.

There are several competing definitions of fairness in machine learning. For example, demographic parity requires similar rates of positive outcomes across different demographic groups, whereas equal opportunity focuses on ensuring that qualified individuals have similar chances of receiving a positive outcome. Other approaches, such as equalized odds, require comparable error rates between groups. These definitions capture different aspects of fairness and can therefore lead to different conclusions about whether the same AI system is fair.

The choice of metric can have a substantial practical impact. Consider an AI system used to select candidates for employment. If demographic parity is used, the system may be designed to select candidates from different demographic groups at similar rates. This can help address historical disadvantages, but it may also produce different decisions than a system focused primarily on equal treatment of similarly qualified candidates.

Consequently, the selection of a fairness metric is also a normative decision. Developers, organisations, policymakers, domain experts, and potentially affected communities may have different views about which inequalities should be reduced and which differences are relevant to the decision-making process. A metric that appears appropriate from a technical perspective may therefore be inappropriate when considering the social context and consequences of the application.

Furthermore, fairness metrics can sometimes be incompatible with one another. Optimising a model for one definition of fairness does not necessarily make it fair according to another definition. This means that there is no universally applicable fairness metric that can simply be implemented in every machine learning system. Instead, the appropriate definition depends on the purpose of the system, the type of decision being made, the potential harms caused by errors, and the values of the stakeholders involved.

\section{Fair AI does not necessarily mean fair outcomes}
Ensuring that an AI system satisfies established fairness criteria does not necessarily guarantee fair outcomes in practice. Fairness can be evaluated at the level of the data, the model, and the fairness target, but the eventual outcome also depends on how the model's predictions are interpreted and used within a broader sociotechnical context. Consequently, unfairness may arise even when the model itself is considered fair.

Human involvement in the decision-making process can introduce unfairness even when the AI system itself produces fair predictions. For example, consider a hiring system that assigns candidates a prospect score based on relevant qualifications and predicts that two candidates have an equal likelihood of succeeding in a position. If the human recruiter treats these identical scores differently. For instance, by accepting the score as sufficient evidence for one candidate while questioning or discounting it for another. The final hiring decision can become unfair despite the underlying model being fair. This may occur when recruiters rely on pre-existing assumptions, stereotypes, or subjective impressions when interpreting the model's output. In this case, the source of unfairness does not lie in the prediction itself, but in the way the prediction is incorporated into the human decision-making process. Similarly a user may decide to investigate a negative AI score for favoured groups, thinking it made a mistake, while they would not do this for an unfavoured group. Thus, equal algorithmic treatment does not necessarily result in equal treatment at the decision-making stage.
\todo{TODO: Mention Robert Williams facial recognition case? Example of improper use, but not necessarily bias. https://quadrangle.michigan.law.umich.edu/issues/winter-2024-2025/flawed-facial-recognition-technology-leads-wrongful-arrest-and-historic}

Finally, unfairness can exist even when the AI system is fair and no additional bias is introduced during the hiring process. The pool of candidates considered by the system may itself be shaped by broader social and economic inequalities. For example, many individuals may have the skills, motivation, and potential to succeed in a particular position but may never reach the minimum qualifications because they have had unequal access to education, training, professional networks, or other opportunities. A hiring system that consistently and fairly selects candidates who meet these requirements may therefore produce fair decisions according to its own criteria while still contributing to an unequal distribution of opportunities. In this case, the source of unfairness lies outside the hiring system itself. This demonstrates that even a completely fair decision-making process cannot, by itself, eliminate unfairness originating from broader structural conditions.

\section{conclusion}
Fairness-aware AI provides valuable tools for reducing specific forms of unfairness in automated decision-making. The results of this thesis contribute to this goal by providing a better understanding of how fairness constraints behave during model training, including the effects of when constraints are applied and how multiple constraints interact. These insights can help practitioners make more informed choices when designing and training fairness-aware models.

At the same time, this reflection highlights that technical fairness should not be equated with fairness in a broader societal sense. A fairness constraint can ensure that a model satisfies a predefined mathematical criterion, but it cannot determine whether that criterion is appropriate for the application or whether the resulting decisions are fair in their broader context. The choice of fairness definition is therefore not merely a technical decision, but also involves normative considerations about which inequalities should be addressed and whose interests should be taken into account.

Furthermore, even a model that satisfies the desired fairness criteria can be incorporated into an unfair decision-making process. Human decision-makers may interpret identical predictions differently, give greater evidential weight to errors affecting certain groups, or use predictions as confirmation of pre-existing beliefs. Beyond the decision-making process itself, structural inequalities may influence who is able to participate in the system in the first place. Therefore, fairness-aware AI should be understood as one component of a broader approach to developing fairer sociotechnical systems. The objective should not be to establish that a model is simply "fair", but to understand which forms of unfairness the model can address, which remain outside its scope, and how the system surrounding the model affects the final outcome. Fairness constraints can make an important contribution by reducing measurable disparities within the model, but they cannot, on their own, resolve the social and structural conditions that produce unfairness. Recognising this limitation does not diminish the value of fairness-aware AI. Rather, it clarifies how such techniques can be used responsibly as part of a wider effort towards fairer decision-making.

